% Auto-generated by paper/build_paper_b.py from paper-B-dynamic-sep-limit.md.
% Edit the markdown source and rerun the build script.
\documentclass[11pt]{article}

\usepackage[T1]{fontenc}
\usepackage[utf8]{inputenc}
\usepackage{lmodern}
\usepackage{microtype}
\usepackage[margin=1in]{geometry}
\usepackage{amsmath,amssymb,bm}
\usepackage{array}
\usepackage{hyperref}
\usepackage{xurl}

\hypersetup{
  colorlinks=true,
  linkcolor=blue,
  urlcolor=blue
}

\title{Dynamic sensitivity as the residual free-fall observable: a first upper limit on lag-responding equivalence-principle violation from the pulsar triple PSR J0337+1715}
\author{Juneyoung, Kim}
\date{2026-07-12}

\begin{document}
\maketitle

\begin{center}
\small
draft manuscript (Paper B of the two-paper split; dynamic free-fall sector)\\
Repository: \texttt{lpaiu-cs/self-mass-unobservability}, branch \texttt{lpaiu/minimal-nonlinear-sideband}\\
Draft 2026-07-12. Companion to Paper A (the static finite-family collapse theorem). Source of record for all numbers: the \texttt{request10\_external/} artifact tree and the \texttt{notes/REQUEST10\_*} request chain on branch \texttt{lpaiu/minimal-nonlinear-sideband} (commit-chain audit in Appendix D).
\end{center}

\begin{abstract}
The companion paper proves that, inside a fixed-order parity-even free-fall worldline EFT domain, every observable consequence of a compact body's self-gravity collapses into finitely many \emph{static} sensitivity coordinates --- coordinates that a pulsar-timing orbital fit absorbs essentially completely. The theorem's boundary map identifies exactly one escape that generates qualitatively new free-fall data at fixed order: an orbital-timescale internal state (boundary \texttt{A4}), whose salvaged bookkeeping is a finite state pair \texttt{(beta, tau\_chi)} --- the amplitude and relaxation lag of a \emph{dynamically responding} sensitivity. This paper instruments that pair directly. We promote the strong-equivalence-principle (SEP) parameter \texttt{Delta} of the pulsar triple PSR J0337+1715 --- statically constrained at the \texttt{2 x 10\textasciicircum{}-6} level and statically 99.94\% degenerate with the timing model --- to a lag-responding transfer sampled at the system's three orbital carriers, and we construct \emph{exact} dynamic response templates by making \texttt{Delta} time-dependent inside the published nonlinear three-body timing integrator. The oscillating-\texttt{Delta} signature proves structurally non-degenerate with the timing model (11-28\% survival per column, versus 0.06\% for static \texttt{Delta}), a property we validate end-to-end against the live integrator, including pulse-number (turn) ambiguity and phase-convention robustness. A pre-registered, gate-audited pipeline on 12,474 public Nancay TOAs (2013-2021) yields no detection and a first upper limit on the amplitude of any lag-responding SEP oscillation:

\begin{verbatim}
|delta Delta| < 1.3 x 10^-9   (95% CL, tau_chi = 2 d, worst drive phase),
rising to 5.0 x 10^-9 at tau_chi = 200 d over the window [2, 500] d,
\end{verbatim}

with a purely statistical floor of \texttt{1.9 x 10\textasciicircum{}-10} and an ultra-conservative bracket of \texttt{1.2 x 10\textasciicircum{}-7} under an unresolved-anchor hypothesis that the data themselves disfavor. In amplitude the dynamic channel probes SEP structure three orders of magnitude below the static bound's absorbable scale and eight orders below the raw static lever arm; in coupling units it is complementary to (not competitive with) the clock-sector bound derived from the same data. The result closes the loop the collapse theorem opened: the only fixed-order free-fall observable the theorem leaves open is now measured, bounded, and reproducible from committed artifacts.
\end{abstract}

\section{Introduction}

Free-fall tests of the strong equivalence principle ask whether a body's gravitational self-energy alters the ratio of its gravitational to inertial mass. The pulsar triple system PSR J0337+1715 provides the cleanest strong-field laboratory: a millisecond pulsar in a 1.63-day inner orbit, with an outer white dwarf in a 327-day orbit supplying the polarizing field. Published analyses bound the static violation parameter \texttt{Delta} (\texttt{m\_g = (1 + Delta) m\_i} for the pulsar) at the few-parts-per-million level.

The companion paper (Paper A) reframes what such static bounds can, even in principle, contain. Inside a fixed-order (\texttt{Delta\_op <= 4}), parity-even, nonspinning, local free-fall worldline EFT domain, the admissible operator space collapses onto finitely many static sensitivity coordinates; a timing-model fit repackages those coordinates into its standard parameters. Concretely, in the present data set the derivative of the residuals with respect to static \texttt{Delta} is enormous --- \texttt{|d res / d Delta| \textasciitilde{} 3.4 x 10\textasciicircum{}10 us} (a \texttt{Delta} of \texttt{1.8 x 10\textasciicircum{}-6} corresponds to 61 ms of raw orbital-polarization signal) --- yet 99.94\% of that column lies inside the span of the 28 fitted timing parameters. The static SEP test is, in the theorem's language, a measurement of the small non-absorbable remainder of one collapsed coordinate.

Paper A's boundary map also identifies where the collapse \emph{stops}: dropping assumption \texttt{A4} (no orbital-timescale internal state) breaks the theorem's Y-only reduction, and the sharp salvage (\texttt{F-A4+}) shows that what survives is not an uncontrolled tower but a finite state datum --- a response amplitude and a relaxation time \texttt{(beta, tau\_chi)} that must be carried separately from all Wilson coefficients. That datum is the \emph{only} qualitatively new free-fall observable the fixed-order theory permits. This paper measures it.

Two facts make the measurement more than a reparameterized static test. First, a \emph{dynamically responding} \texttt{Delta(t)} --- the sensitivity lag-tracking the driving potential through a one-state transfer --- produces timing structure (growing orbital-phase modulation and carrier-quadrature shifts) that the timing model cannot reproduce: the per-column nuisance survival is 11-28\%, versus 0.06\% for the static column. The degeneracy that makes the static test hard is \emph{absent} for the dynamic signature. Second, the response can be computed exactly: \texttt{Delta} is implemented natively in the published nonlinear three-body integrator, and a ten-line, integrity-gated extension makes it time-dependent, turning the integrator itself into the template generator --- no adiabatic or secular approximation enters the analysis.

The remainder of the paper: Section 2 defines the observable and its carrier/comparator structure; Section 3 constructs the exact templates; Section 4 specifies the pre-registered estimation pipeline and its anchor brackets; Section 5 reports the validation-gate record, including two measured findings of independent interest (the turn-fixing behavior of the integrator interface, and a quantification of chi2-free gap-hidden pulse-number slips); Section 6 gives the results; Sections 7 and 8 state scope and outlook.

\section{The dynamic sensitivity observable}

\subsection{One-state transfer}

The minimal \texttt{A4}-escape model attaches to the pulsar one internal state \texttt{chi} with linear relaxation dynamics driven by an ambient-field invariant \texttt{F(t)}:

\begin{verbatim}
tau_chi  d chi/dt + chi = alpha F(t),
Delta_eff(t) = Delta_static + c_Y F(t) + c_chi chi(t).
\end{verbatim}

For a drive component at angular frequency \texttt{w}, the readout transfer is

\begin{verbatim}
G(i w) = c_Y + beta / (1 + i w tau_chi),        beta = alpha c_chi :
\end{verbatim}

an instantaneous term plus a single-pole lag. The pole is the payload: its frequency dependence across distinct carriers is what no finite instantaneous (derivative-polynomial) comparator can reproduce.

\subsection{Carriers and comparator pressure}

The triple system supplies three nonresonant drive carriers,

\begin{verbatim}
Omega_in = 3.856137 rad/d,  Omega_out = 0.019200 rad/d,
Omega_dif = |Omega_in - Omega_out| = 3.836937 rad/d
(ratio 200.8; 1:1, 2:1, 1:2 excluded),
\end{verbatim}

and the request-chain counting theorems fix what three samples of \texttt{G} can distinguish: a real shared-coefficient derivative comparator is pressured through degree 4, and the deliberately generous complex-coefficient polynomial comparator through degree 1 (\texttt{K\_required = N + 2}). The carrier projection is not arbitrary per-carrier freedom: the physical projection manifold is phase-locked (five real parameters in the six-dimensional carrier space), and the timing model's finite span realizes precisely such a manifold --- measured rank \texttt{<= 5/6} at every threshold, unchanged when the planet-block Jacobian columns are re-derived at perturbative steps.

\subsection{Free-fall implementation}

In the integrator of record (Nutimo; Zenodo 13899771), \texttt{Delta} enters as a rescaling of the pulsar's pairwise gravitational coupling, \texttt{G\_eff(pulsar, j) = (1 + Delta) G}, applied symmetrically in the accelerations of the pulsar and of both companions --- the full nonlinear three-body-plus-planet polarization dynamics of the published static test. The dynamic observable promotes this coupling:

\begin{verbatim}
Delta(t) = Delta_static
         + Re sum_w [ (c_Y + beta/(1 + i w tau_chi)) d_w e^{i(w t + phi_w)} ],
\end{verbatim}

with unit drive (\texttt{d\_w = 1}) as the convention-clean primary normalization (Section 6.3 gives the dictionary-anchored coupling version). Because \texttt{Delta} itself is the dimensionless SEP parameter, an amplitude bound on \texttt{delta Delta} is a physical statement independent of the drive model.

\section{Exact response templates from the nonlinear integrator}

\subsection{The time-dependent-Delta extension}

A surgical patch (seven count-asserted edits, one file) makes the coupling time-dependent at its unique active application site,

\begin{verbatim}
G_eff(i, j; t) = Gg[i][j] + A cos(w t + phi)   on pulsar pairs only,
\end{verbatim}

read from the environment at each initialization, built in a separate source tree. The integrity stop-rule holds exactly: at \texttt{A = 0} the patched build reproduces the baseline residuals with \texttt{max |diff| = 0.0 us} (bit-exact).

\subsection{Time-convention calibration and the integral-kernel signature}

The integrator's internal time unit was calibrated empirically after two line-search designs were stopped by their own pre-registered rules (the response to an oscillating coupling is broadband-dominated, defeating single-line and naive-spectral tests). The decisive probe is quasi-static: a ramp drive (period 30,000 d, phase pi/2) fitted against the \emph{measured} static column. Both unit hypotheses fit with \texttt{R\textasciicircum{}2 = 0.999} (deep-adiabatic shape degeneracy) but the amplitudes decide, and they decide coherently: \texttt{a\_hat(timedays) = 0.502} and \texttt{a\_hat(days) = 0.0219}, exactly the \texttt{1/2} and \texttt{1/(2 x 24.077)} predicted if (i) the integrator time unit is 24.0774 days and (ii) the secular-dominated response accumulates as the \emph{integral} of \texttt{Delta(t)} (a phase-drift kernel), for which a ramp yields half the instantaneous-factorization template. One measurement therefore fixes the convention and validates the patch normalization end-to-end.

\subsection{Production columns}

Six exact response columns were derived --- cosine and sine drive quadratures at each carrier --- by central differences in \texttt{A} within the verified linear regime (adaptive amplitudes down to \texttt{3.3 x 10\textasciicircum{}-10}, where the inner-carrier response reaches \texttt{2.6 x 10\textasciicircum{}12 us} per unit oscillation amplitude; the turn-wrap onset at \texttt{|Delta| \textasciitilde{} 10\textasciicircum{}-7} bounds the usable regime and was mapped). Any \texttt{(c\_Y, beta, tau\_chi, phase)} template is a closed-form linear combination of the six columns; no kernel or adiabatic assumption enters.

The templates' key structural property, measured against the full nuisance span of Section 4: per-column survival 11-28\%, shared-\texttt{tau} template survival peaking at 43.8\% (\texttt{tau = 17.7 d}) --- three orders of magnitude less absorbable than the static column (0.06\%). The timing model cannot imitate an oscillating coupling.

\section{Pre-registered estimation pipeline}

\subsection{Estimator}

All quoted statistics derive from a deterministic, artifacts-only joint weighted fit (Frisch-Waugh-Lovell block elimination; seed 20260710):

\begin{itemize}
\item \textbf{Data}: 12,474 public Nancay TOAs (2013-2021, span 2987.9 d), baseline residuals of the published maximum-posterior solution (wrms 1.887 us, EFAC 1.1225).
\item \textbf{Nuisance block}: the 28 fitted-parameter response columns (with the seven planet-block columns re-derived at perturbative steps after the original one-sigma-scale steps were shown to be nonperturbative secants), an offset column, thirty low-frequency Fourier pairs (prior-free red-noise marginalization), and the static-\texttt{Delta} column (static-leakage guard); truncated at relative singular value \texttt{10\textasciicircum{}-3}. Noise scale \texttt{s = 1.115} (\texttt{chi2\_red = 1.243}).
\item \textbf{Signal block}: \texttt{[T\_cY, T\_beta(tau)]} from the measured columns; \texttt{tau} grid of 62 points on \texttt{[2, 500] d}.
\item \textbf{Detection}: \texttt{Z = max |z|} over the full grid, calibrated by 500 noise simulations through the identical (trials-inclusive) statistic; an anti-causal control (the unphysical advanced-response template, \texttt{g2 -> -g2}) is scanned identically and must stay quiet.
\item \textbf{Limits}: flat-prior Gaussian-posterior 95\% amplitudes \texttt{u95(tau)}.
\item \textbf{Phase marginalization}: the one unresolved convention is a common time-origin \texttt{t\_off}; carrier phases rotate as \texttt{w t\_off} (independent per-carrier phases belong to the excluded arbitrary-projection collapse class). The quoted limit is the \emph{worst phase} per \texttt{tau} over 24 origins spanning the inner period --- valid whatever the true phase. The marginalization costs only 12\%, confirming phase-robust non-degeneracy.
\end{itemize}

\subsection{Anchor brackets}

The Fisher machinery's absolute scale was audited rather than assumed. A ratio spectrum \texttt{rho\_j = sigma\_MCMC,j / sigma\_Fisher,j} over the twenty core timing parameters (using the released chain widths) rejects global Fisher optimism: \texttt{rho <= 1} everywhere, with \texttt{rho = 0.8-0.94} precisely where the comparison is convention-clean. The \texttt{\textasciitilde{}10\textasciicircum{}3} gap between the Fisher and published widths of the \emph{static} \texttt{Delta} direction is therefore direction-specific --- the turn-wrap manifold that inflates the static posterior demonstrably does not act on the dynamic templates (Section 5). Per a pre-registered rule the working anchor is \texttt{K\_dyn = 10} (a safety floor, not a measurement), and every quoted quantity carries the full bracket: the statistical Fisher floor, the \texttt{K\_dyn = 10} headline, and the \texttt{K = 934} ultra-conservative fallback.

\section{Validation-gate record}

Every stage ran under pre-registered, committed rules; each amendment was itself committed before the stage it governed. The complete record:

\begin{center}
\begin{tabular}{p{0.07\textwidth}p{0.14\textwidth}p{0.76\textwidth}}
\hline
Gate & Question & Result \\
\hline
F0 & is the static column a true derivative? & half-step linearity \texttt{0.0000} at \texttt{h = 10\textasciicircum{}-8}; wrap onset mapped at \texttt{\textasciitilde{}10\textasciicircum{}-7} \\
A=0 & does the patched integrator reproduce baseline? & \texttt{max diff = 0.0 us} (bit-exact) \\
Ramp & time convention + normalization & integral-kernel \texttt{1/2} signature, both hypotheses coherent (Sec. 3.2) \\
F2 & dynamic non-degeneracy & survival 11-28\% per column; 43.8\% shared-\texttt{tau} peak \\
F3 & projected sensitivity & anchored \texttt{6 x 10\textasciicircum{}-8 <= 10\textasciicircum{}-5} threshold (x160 margin) \\
G1 & live null calibration & one Gauss-Newton fit of the real residuals with live integrator evaluations; ` \\
G2a & estimator calibration, detection regime & \texttt{4 sigma} injections recovered within \texttt{0.06 sigma} (undamped live GN) \\
G2b & estimator calibration, limit amplitude & injections at the conservative limit (\textasciitilde{}48,000 whitened units) recovered to 0.16-0.23\% relative through step-capped live GN; the live model absorbs essentially nothing \\
Turn V & does the interface re-assign pulse numbers? & no: alias-shift response exceeds \texttt{P/2} (1522 vs 1366 us) --- turns fixed, wrap arithmetic definitional \\
Turn L & can any turn-aliased solution mimic/hide the template? & 14 chi2-viable alternative turn solutions exist per amplitude (integer slips hidden in observing gaps --- the phase-connection ambiguity, quantified against the full span), and the estimator is stable across \emph{all} of them: worst \texttt{beta} deviation \texttt{5.6 x 10\textasciicircum{}-11} vs \texttt{1.8 x 10\textasciicircum{}-10} tolerance \\
\hline
\end{tabular}
\end{center}

Two by-products deserve independent note. First, the turn-fixing behavior of the interface is now a measured fact, not a documentation claim. Second, the lattice scan is, to our knowledge, the first quantification of the gap-hidden pulse-number-slip degeneracy against a full timing-model span for this data set --- and simultaneously a proof that the degeneracy is harmless for the dynamic observable: pulse-number freedom is chi2-free but orthogonal to the template in the estimator's measurement space.

\section{Results}

\subsection{No detection}

Across every pipeline mode --- zero-phase, phase-marginalized (trials-corrected \texttt{Z = 1.91}, global \texttt{p = 0.30}), and dictionary-anchored (\texttt{p = 0.53}) --- the carriers are noise-consistent, and the anti-causal control is quiet (\texttt{p = 0.27}). Linear injection-recovery is exact at all anchors; live-model injection-recovery is exact to sub-sigma (Section 5).

\subsection{The upper limit}

\begin{verbatim}
95% CL upper limit on the amplitude of a lag-responding SEP oscillation,
phase-marginalized (worst phase), K_dyn = 10, window tau_chi in [2, 500] d:

  tau_chi        2 d       5 d      18 d      52 d      200 d
  |dDelta| <   1.30e-9   1.36e-9  1.42e-9   1.84e-9   4.95e-9

  statistical (Fisher) floor:     1.9e-10  (tau = 2 d)
  ultra-conservative (K = 934):   1.2e-7   (tau = 2 d)
  zero-phase convention curve:    1.17e-9  (tau = 2 d; recorded)
\end{verbatim}

Interpretation: any mechanism --- an internal mode, a light dynamical degree of freedom, a relaxation response of the pulsar's effective self-gravity --- that makes the SEP parameter of this neutron star \emph{oscillate} at the triple system's carriers with lags between 2 and 500 days is excluded at amplitudes three orders of magnitude below the static bound, and the exclusion is robust to the timing model, pulse-number ambiguity, drive phase, and noise model at the levels recorded in Section 5.

\subsection{Coupling units and channel complementarity}

Under the leading-order dictionary drive (\texttt{F = U/c\textasciicircum{}2} at the pulsar; dimensionless carrier amplitudes \texttt{f\_w \textasciitilde{} 10\textasciicircum{}-11 - 10\textasciicircum{}-10} with tasc-locked phases; O(1) geometric factors set to unity), the same fit gives the Conjectural coupling bound \texttt{beta\_phys < 14-17} for \texttt{tau in [2, 52] d}. The clock-sector analysis of the same data (the 10.7 chain) bounds its rate-coupling composite at \texttt{0.4-0.6}: the free-fall drive lacks the clock channel's \texttt{1/Omega} integration amplification, so in coupling units the two channels bound \emph{different composites} of the underlying state data, while in raw \texttt{Delta} amplitude the free-fall channel is deeper by \texttt{3 x 10\textasciicircum{}8}. The two bounds are complementary faces of the same \texttt{(beta, tau\_chi)} datum.

\section{Scope: what is and is not claimed}

Claimed: a first, validated, reproducible 95\% upper limit on the amplitude of lag-responding SEP violation of PSR J0337+1715 over \texttt{tau\_chi in [2, 500] d}, in the stated conventions, with the full anchor bracket published.

Not claimed: any detection; any statement below \texttt{tau\_chi = 2 d} (the \texttt{c\_Y}-degeneracy grows toward the instantaneous limit; sub-edge rows are diagnostic); any statement for other systems or other state models (the one-pole transfer is the minimal \texttt{F-A4+} branch; higher-order state dynamics would need more carriers per the counting boundaries); an absolute maximum-likelihood red-noise treatment (the prior-free Fourier marginalization shifted the clock-channel limits by only 6-10\%); physical drive factors beyond leading order.

\section{Discussion}

The measurement completes a loop that began as a slogan. Paper A converts "matter cannot sense its own mass" into a conditional collapse theorem whose sharp boundary map does two things at once: it explains why static free-fall SEP structure is largely invisible (absorbed into the fit), and it points to the unique fixed-order escape carrying genuinely new data. This paper builds the instrument for that escape --- exact dynamic templates from the published integrator itself --- and returns the first number: \texttt{|delta Delta| < 1.3 x 10\textasciicircum{}-9}. The method (pre-registration with committed amendments, live-model gates, adversarial controls, bracketed anchors) travelled two failed template-calibration designs, one integrator segfault, one refuted absorber hypothesis, and one mis-scaled gate without ever quoting an unvalidated number; we consider that audit trail as much a result as the limit.

Three continuations are natural. A second, independent data set (the published multi-telescope TOAs) would test the limit's pipeline-dependence and could tighten the anchor bracket from the data side. The sub-2-day window --- where the pole approaches instantaneity --- needs the two-quadrature degeneracy treated as a joint confidence region rather than a profiled bound. And the anchor question itself (\texttt{K\_dyn}) reduces to reproducing the published static posterior within a linearized-plus-turn-manifold model; the lattice machinery of Section 5 is the starting point.

\section{Appendix A: request-chain map}

\texttt{notes/REQUEST10\_DYNAMIC\_CHI\_MULTIFREQUENCY.md} (10.1, transfer/comparator counting) -> \texttt{REQUEST10\_TRIPLE\_SHARED\_TAU\_CARRIER\_BRIDGE.md} (10.2b) -> \texttt{REQUEST10\_TRIPLE\_GR\_CARRIER\_INVENTORY.md} (10.3, three-carrier boundary) -> \texttt{REQUEST10\_TRIPLE\_PROJECTION\_NUISANCE\_GATE.md} / \texttt{REQUEST10\_PHYSICAL\_PROJECTION\_MANIFOLD\_GATE.md} (10.4/10.5, projection classes) -> \texttt{REQUEST10\_NAMED\_TIMING\_MODEL\_PROJECTION\_AUDIT.md} (10.6) -> \texttt{REQUEST10\_EXTERNAL\_NUTIMO\_HANDOFF\_PACKET.md} + \texttt{request10\_external/} (10.7a pilot) -> \texttt{REQUEST10\_EXTERNAL\_JOINT\_FIT\_UPPER\_LIMIT.md} (10.7b) -> \texttt{...\_AMENDMENT\_10\_7C.md} -> \texttt{...\_PROMOTION\_10\_7D.md} -> \texttt{...\_PHYSICAL\_ANCHOR\_REDNOISE\_10\_7E.md} (clock-channel chain) -> \texttt{REQUEST10\_8\_DYNAMIC\_SEP\_POLARIZATION\_CHANNEL.md} (10.8 design + F-gates) -> \texttt{REQUEST10\_8B\_REALDATA\_DYNAMIC\_SEP\_FIT.md} -> \texttt{REQUEST10\_8C\_ANCHOR\_RESOLUTION.md} -> \texttt{REQUEST10\_8D\_TURN\_SEARCH\_ROBUSTNESS.md} -> \texttt{REQUEST10\_8E\_PHASE\_MARG\_PHYSICAL\_ANCHOR.md} -> \texttt{REQUEST10\_8\_REVIEW.md}.

\section{Appendix B: artifacts and reproduction}

All quoted numbers reproduce deterministically from \texttt{request10\_external/} with stock numpy: the quoted curves from \texttt{scripts/sep\_phase\_marg\_10\_8e.py} (headline) and \texttt{scripts/sep\_joint\_fit\_10\_8b.py} (zero-phase), gates from \texttt{scripts/sep\_feasibility\_gates.py} and \texttt{scripts/turn\_search\_10\_8d.py}, anchor diagnosis from \texttt{sep\_dynamic/anchor\_resolution\_10\_8c.json}. Inputs of record: \texttt{baseline\_planetGR.npz}, \texttt{finite\_jacobian\_v2.npy} (+ \texttt{jac\_v2/}), \texttt{sep\_dynamic/col\_SEP\_D.npz}, \texttt{sep\_dynamic/sep\_dynamic\_columns.npz}, and the integrator patch \texttt{sep\_dynamic/patches/sepdyn.patch}. Live-model gates ran in the WSL runtime described in the 10.7a packet (Zenodo 13899771 build).

\section{Appendix C: gate-number table}

F0: \texttt{h = 10\textasciicircum{}-8}, maxdev 340 us, half-step dev 0.0000. A=0: 0.0 us. Ramp: \texttt{a\_hat = 0.502 / 0.0219} (predicted 0.5 / 0.0208). F2: best survival 0.438 at \texttt{tau = 17.7 d}. F3: anchored \texttt{6.1 x 10\textasciicircum{}-8}. 10.8b fit: \texttt{Z = 0.318}, \texttt{p = 0.95}, anti-causal \texttt{p = 0.48}, \texttt{s = 1.1148}, nuisance rank 70/90. G1: \texttt{|dz| <= 0.019}. G2a: dev \texttt{<= 0.062 sigma\_F}. G2b: rel. dev \texttt{<= 0.23\%}. Turn V: 1522 us > P/2 = 1366 us. Turn L: 14 viable cells; worst \texttt{beta} dev \texttt{5.6 x 10\textasciicircum{}-11} (tol \texttt{1.8 x 10\textasciicircum{}-10}). 10.8c: core \texttt{rho} max 0.84. 10.8e: trials-corrected \texttt{Z = 1.91 (p = 0.30)}; worst-phase cost 12\%.

\section{Appendix D: commit-chain audit (pre-registrations and amendments)}

Clock chain: \texttt{3d9337c} (10.7b pre-reg) -> \texttt{92667cf} (10.7c) -> \texttt{a88d27f} -> \texttt{e7bf8e9} (10.7d) -> \texttt{43a6618} -> \texttt{e7f5f5a} (10.7e) -> \texttt{aa50c9b}. Free-fall chain: \texttt{ac13406} (10.8 design) -> \texttt{69dd1c1} (F0/F1) -> \texttt{fa1fb02} (R5 refuted) -> \texttt{2515481} (patch + A=0) -> \texttt{b3633cf} (T2/F2/F3) -> \texttt{2282924} (10.8b pre-reg) -> \texttt{6e55569} (G2 amendment) -> \texttt{426ad2c} (10.8b) -> \texttt{26abfa1} (10.8c pre-reg) -> \texttt{99f0ab3} (10.8c) -> \texttt{349c70c} (10.8d pre-reg) -> \texttt{5df4870}/\texttt{94c04c3} (amendments) -> \texttt{14c950d} (10.8d) -> \texttt{3827649} (10.8e pre-reg) -> \texttt{3556363} (10.8e). Every amendment precedes the stage it governs.

\section{References (draft)}

\begin{enumerate}
\item S. M. Ransom et al., "A millisecond pulsar in a stellar triple system," Nature 505, 520 (2014).
\item A. M. Archibald et al., "Universality of free fall from the orbital motion of a pulsar in a stellar triple system," Nature 559, 73 (2018).
\item G. Voisin et al., "An improved test of the strong equivalence principle with the pulsar in a triple star system," A\&A 638, A24 (2020).
\item G. Voisin et al., data and code release, Zenodo 13899771 (Nutimo; A\&A 2024, doi:10.1051/0004-6361/202452100).
\item T. Damour, "Binary systems as test-beds of gravity theories," in Physics of Relativistic Objects in Compact Binaries (2007) --- the \texttt{(Delta, gammabar, betabar)} strong-field parameterization.
\item C. M. Will, "The confrontation between general relativity and experiment," Living Rev. Relativity 17, 4 (2014).
\item W. D. Goldberger and I. Z. Rothstein, "An effective field theory of gravity for extended objects," Phys. Rev. D 73, 104029 (2006).
\item K. Nordtvedt, "Equivalence principle for massive bodies," Phys. Rev. 169, 1014 (1968).
\item J. Kim, "Finite-Family Collapse of Free-Fall Self-Energy Couplings: A Fixed-Order Worldline-EFT Theorem, its Uniqueness No-Go, and Sharp Boundary Escapes" (Paper A, companion), this repository, \texttt{paper/paper-A-collapse-theorem.md} (branch \texttt{main}).
\end{enumerate}

\end{document}
